\chapter{多巴胺到底在说什么}
% 完整版：chapters/07_dofaalt.tex

在上一章里，多巴胺很简单：一个数，“比预期的好还是差”。这是一阶近似，而且很管用。但近些年的实验表明，这根“导线”里传的不止一个数。本章讲的是图景更复杂的地方，以及科学家仍在争论的地方。

\section*{乐观者与悲观者}

多巴胺神经元并不一样。有些对惊喜的反应比对坏意外更强——这是乐观者。另一些正好相反——是悲观者。如果每个神经元按自己的规则学习，乐观者会学到最好情况下奖励有多少，悲观者则学到最坏情况下有多少。它们合在一起记下的不是一个平均预期，而是所有可能结果的整个分布，就像一个博彩庄家，不仅知道谁是热门，还知道每一方的赔率。这个想法与实验相符；至于大脑为什么需要整个分布——目前还是假说。

\pencilfig{7}{\pencilart{7}}{每个多巴胺神经元守住可能奖励曲线上的一个点：合起来，它们记住了整条曲线。}

\section*{不只是“更好”，还有“重要”}

有一部分多巴胺神经元对不愉快但重要的事件也会爆发：巨响、危险。看来，多巴胺网络传递的不只是意外的符号，还有它的显著性——一个“注意！”的信号。

\section*{一根导线上的两个信号}

快速的爆发负责教学。而多巴胺缓慢的平均水平以分钟为单位变化，说的似乎是另一回事：眼下的环境有多富足。如果周围奖励很多，时间就很宝贵——动物会行动得更快。实验证实，多巴胺水平与动作的干劲有关。工程师会把这叫作按频率分离信号：高频携带一种信息，低频携带另一种。

\section*{接近目标时的爬升}

当动物接近奖励时，多巴胺常常不是爆发，而是平稳上升。一种解释是：这就是预期本身，一个“目标近了，加把劲”的信号。另一种解释保留了意外的思路：动物离得远时对环境的估计很模糊，越走近越精确，而每一次修正都是一个小小的惊喜。有一个巧妙的实验支持后者：在虚拟走廊里，把动物瞬间挪到离目标更近的地方，多巴胺便以一次爆发作答，正如意外应有的那样。

\section*{回头看}

还有一种大胆的替代理论：多巴胺回答的不是“将会发生什么？”，而是“是什么带来了奖励？”——它看向过去，而不是未来。在两种理论预测不同的实验里，结果目前还互相矛盾。这是一场正在进行的科学争论。

\section*{向量而非数值}

最后，当奖励的价值不变、只是\emph{口味}变了时，多巴胺也会有反应，而且它还能帮助在完全没有奖励的情况下学习。看来它不是一个误差，而是一整束误差——所发生事情的每一个特征各对应一个。

本章的结论：大多数新理论并没有推翻那幅简单的图景，而是扩展了它。多巴胺不是一根导线，而是由几根导线组成的一小束线缆。只有“回头看”这一理论真正向它发起了挑战。

\fullversion{7}
